% =====================================================================
%  02_marco-teorico.tex — Conceptos y tecnologías en que se apoya SARU
% =====================================================================
\chapter{Marco teórico}
\label{cap:marco}

Este capítulo resume los conceptos en que se apoya \saru{}: la detección y
el seguimiento de objetos en video, la lectura de caracteres, la
determinación de cruces con una línea virtual, la computación en el borde y
las tecnologías de soporte del sistema.

\section{Visión artificial y detección de objetos}

La detección de objetos consiste en localizar, dentro de una imagen, cada
instancia de una clase de interés y delimitarla con un recuadro. Los
detectores modernos se basan en redes neuronales convolucionales y se
dividen en dos familias: los de dos etapas, que primero proponen regiones y
después las clasifican, y los de una etapa, que predicen recuadros y clases
en una sola pasada.

\subsection{YOLO}

YOLO (\emph{You Only Look Once}) es un detector de una etapa que divide la
imagen en una cuadrícula y predice, para cada celda, recuadros con su
confianza y su clase en una sola evaluación de la red~\cite{redmon2016yolo}.
Al no requerir una etapa de propuesta de regiones, es lo bastante rápido
para procesar video en tiempo real. \saru{} usa la implementación de
Ultralytics~\cite{ultralytics} para sus tres modelos.

Cada detección tiene una \textbf{confianza} entre 0 y 1; el sistema
descarta las que quedan por debajo de un umbral. Un umbral alto reduce los
falsos positivos a cambio de perder detecciones dudosas; uno bajo hace lo
contrario.

\subsection{Modelos en cascada}

En lugar de un solo modelo que lo resuelva todo, \saru{} encadena tres
detectores especializados. Cada etapa trabaja sobre el recorte que entrega
la anterior, de modo que los modelos posteriores procesan imágenes más
pequeñas y con menos distractores:

\begin{description}[style=nextline]
  \item[M1 — unidad] Detecta autobuses en el cuadro completo.
  \item[M2 — zona del número] Dentro del recorte de cada unidad, localiza
        la zona donde está pintado el número y la clasifica como
        delantera, lateral o trasera.
  \item[M3 — dígitos] Dentro de la zona, detecta cada dígito del 0 al 9
        como una clase independiente.
\end{description}

\section{Seguimiento multiobjeto}

Detectar una unidad en un cuadro no basta: hay que saber que la unidad del
cuadro siguiente es la misma. El enfoque de \emph{seguimiento por
detección} asocia las detecciones de cuadros consecutivos y asigna a cada
objeto un identificador que se conserva mientras permanece en escena.

ByteTrack es un algoritmo de este tipo que asocia todas las detecciones,
incluidas las de baja confianza, en lugar de descartarlas. Con ello
recupera objetos parcialmente ocultos y reduce los cambios de
identificador~\cite{zhang2022bytetrack}. En \saru{}, el identificador de
seguimiento permite acumular varias lecturas del número de una misma
unidad y evaluar su trayectoria respecto de la línea virtual.

\section{Lectura de caracteres por detección}

Los sistemas de reconocimiento automático de matrículas suelen dividirse
en detección del vehículo, localización de la placa y reconocimiento de
caracteres~\cite{shashirangana2021alpr}. Una alternativa para la última
etapa es tratar cada carácter como un objeto: un detector localiza los
dígitos y estos se ordenan por su posición horizontal para formar el
número~\cite{laroca2018yolo}. \saru{} aplica este enfoque con M3.

Antes de leer los dígitos, la imagen de la zona se mejora con CLAHE
(\emph{Contrast Limited Adaptive Histogram Equalization}), que ecualiza el
contraste por regiones pequeñas y limita la amplificación del
ruido~\cite{zuiderveld1994clahe}. Esto ayuda con números pintados en
superficies con reflejos o sombras.

\subsection{Confirmación por votación}

Una sola lectura puede confundir dígitos parecidos (un 5 con un 6, un 8 con
un 3) o perder uno. Por eso un número solo se \textbf{confirma} cuando
varias lecturas de la misma unidad coinciden. Si un número leído es otro
con dígitos de menos (por ejemplo, \texttt{504} dentro de \texttt{5014}), se
prefiere el más largo.

\section{Determinación de cruces con una línea virtual}

Una \textbf{línea virtual} es un trazo definido por puntos en la imagen de
la cámara. Para saber de qué lado de un segmento $\overline{PQ}$ se
encuentra un punto $B$ de la unidad, se usa el signo del producto cruz:

\begin{equation}
  s = (Q_x - P_x)(B_y - P_y) - (Q_y - P_y)(B_x - P_x)
\end{equation}

Un valor positivo indica un lado y uno negativo el otro. Un cambio de signo
entre dos posiciones consecutivas indica un posible cruce. Para que el
conteo sea robusto se agregan tres reglas:

\begin{itemize}
  \item \textbf{Margen}: dentro de una franja alrededor de la línea el lado
        no cambia. Así, el temblor de la detección con la unidad detenida
        sobre la línea no produce cruces en cadena.
  \item \textbf{Cruce real}: con una línea finita, el cambio de lado solo
        cuenta si el segmento recorrido por la unidad corta algún tramo de
        la línea; rodearla por un extremo no es un cruce~\cite{rep-tecnico}.
  \item \textbf{Maniobra y repetición}: si la unidad regresa antes de cierto
        tiempo, el movimiento se descarta; y un mismo número no se registra
        dos veces en un intervalo corto.
\end{itemize}

\section{Computación en el borde}

La computación en el borde procesa los datos cerca de donde se generan, en
lugar de enviarlos a un centro de datos. Reduce la latencia, el consumo de
ancho de banda y la exposición de datos sensibles~\cite{shi2016edge}. Para
video, evita transmitir de forma continua un flujo de alta resolución.

\subsection{NVIDIA Jetson Orin y TensorRT}

La familia Jetson Orin integra CPU y GPU en un módulo de bajo consumo
pensado para inferencia en el borde~\cite{jetson}. TensorRT optimiza un
modelo entrenado para una GPU concreta: fusiona capas, reduce la precisión
numérica cuando es seguro hacerlo y genera un \emph{motor}
(\texttt{.engine})~\cite{tensorrt}. Ese motor depende de la GPU y de la
versión de TensorRT con que se creó, por lo que debe generarse en la propia
Jetson.

\section{Validación con modelos de lenguaje multimodales}

Los modelos de lenguaje multimodales reciben imágenes y texto y responden
en lenguaje natural. \saru{} usa Gemini~\cite{gemini} como segunda opinión:
le envía la foto del cruce y el número leído, y el modelo indica si
coinciden. Cuando hay duda, la decisión pasa a una persona por Telegram.

\section{Tecnologías de soporte}

\begin{longtable}{@{}L{3.4cm}L{10.6cm}@{}}
  \caption{Tecnologías de soporte de \saru{}.}
  \label{tab:tecnologias} \\
  \toprule
  \textbf{Tecnología} & \textbf{Uso en el sistema} \\
  \midrule
  \endfirsthead
  \toprule
  \textbf{Tecnología} & \textbf{Uso en el sistema} \\
  \midrule
  \endhead
  OpenCV & Lectura del video RTSP y manejo de imágenes~\cite{opencv}. \\
  Flask & Servidor web del tablero y de la API de consulta~\cite{flask}. \\
  SQLite & Base de datos local de cruces y estancias, en modo WAL para
    permitir lecturas mientras se escribe~\cite{sqlite}. \\
  Docker Compose & Despliegue de los tres servicios del sistema en la
    Jetson. \\
  Telegram Bot API & Avisos al grupo de operación y validación con botones
    \cite{telegram}. \\
  Baileys & Cliente de WhatsApp Web para reenviar los avisos; no es la API
    oficial de WhatsApp~\cite{baileys}. \\
  FFmpeg & Grabación continua del video de la cámara sin recodificar. \\
  \bottomrule
\end{longtable}
